How do security scanners compare when evaluated on real-world code? We
present \textsc{RealVuln}, the first open-source benchmark that
systematically compares three categories of vulnerability scanners---rule-based
SAST, general-purpose LLMs, and security-specialized systems---across 26
real-world Python repositories containing 796 hand-labeled entries (676 vulnerabilities and
120 false-positive traps). We evaluate 15 scanners (3 rule-based SAST
tools, 10 general-purpose LLM configurations, and 2 security-specialized
systems) using the recall-weighted F3 score ($\beta{=}3$, recall weighted 9$\times$ over precision) as the primary metric. Our
results reveal a clear three-tier hierarchy: the security-specialized
scanner (Kolega, F3\,=\,73.0) substantially outperforms the best
general-purpose LLM (Claude Sonnet~4.6, F3\,=\,51.7), which in turn
outperforms the best rule-based SAST tool (Semgrep, F3\,=\,17.7) by
nearly 3$\times$. All code,
ground-truth data, scanner outputs, and scoring scripts are released
under an open-source license. An interactive dashboard is available at \url{https://realvuln.kolega.dev/}. \textsc{RealVuln} is designed as a living
benchmark---versioned, community-driven, and with a roadmap toward
multi-language coverage.
